\documentclass[11pt,a4paper]{article}

\usepackage[utf8]{inputenc}
\usepackage[T1]{fontenc}
\usepackage[ngerman]{babel}

\usepackage{amsmath,amssymb,amsthm}
\usepackage{mathtools}
\usepackage[a4paper,margin=2.5cm]{geometry}
\usepackage{lmodern}
\usepackage{microtype}
\usepackage{booktabs}
\usepackage{longtable}
\usepackage{array}
\usepackage{hyperref}
\usepackage{xcolor}

% Robuste Grafikeinbindung
\usepackage{graphicx}
\DeclareGraphicsExtensions{.pdf,.png,.jpg}
\graphicspath{{./}{fig/}{figs/}{images/}{img/}}
\newcommand{\includegraphicssafe}[2][]{%
	\IfFileExists{#2}{\includegraphics[#1]{#2}}{%
		\fbox{\parbox[c][4cm][c]{0.8\linewidth}{\centering Grafikdatei \texttt{#2} nicht gefunden}}%
	}%
}

% Theorem-Umgebungen
\theoremstyle{plain}
\newtheorem{satz}{Satz}
\newtheorem{lemma}{Lemma}
\newtheorem{korollar}{Korollar}

\theoremstyle{definition}
\newtheorem{definition}{Definition}
\newtheorem{beispiel}{Beispiel}

\theoremstyle{remark}
\newtheorem{bemerkung}{Bemerkung}

\title{\textbf{Tabellen zu Zwillingen, Drillingen und Vierlingen\\
		mit Heimatfamilien und Familienindizes}}
\author{Kollaborative Ausarbeitung}
\date{\today}

\begin{document}
	
	\maketitle
	
	\section{Vorbemerkung}
	
	Wir verwenden die mod-\(12\)-Klassifikation
	\[
	E \equiv 1 \pmod{12},\qquad
	A \equiv 5 \pmod{12},\qquad
	B \equiv 7 \pmod{12},\qquad
	C \equiv 11 \pmod{12}.
	\]
	
	Für eine Primzahl \(q>3\) bezeichne \(F_i\) den Index von \(q\) innerhalb ihrer Heimatfamilie \(F\in\{E,A,B,C\}\). Beispielsweise gilt:
	\[
	5=A_1,\qquad 7=B_1,\qquad 11=C_1,\qquad 13=E_1,
	\]
	\[
	17=A_2,\qquad 19=B_2,\qquad 23=C_2,\qquad 37=E_2.
	\]
	
	Die Primzahl \(3\) wird gesondert notiert, da sie nicht zu den vier Familien \(E,A,B,C\) gehört.
	
	\section{Tabelle der ersten Zwillinge}
	
	\begin{longtable}{>{\raggedleft\arraybackslash}p{1.2cm} >{\centering\arraybackslash}p{4.0cm} p{8.2cm}}
		\toprule
		Nr. & Zwilling & Familienindizes \\
		\midrule
		\endfirsthead
		\toprule
		Nr. & Zwilling & Familienindizes \\
		\midrule
		\endhead
		\bottomrule
		\endfoot
		
		1  & $(3,5)$       & $3=\text{Sonderfall},\ 5=A_{1}$ \\
		2  & $(5,7)$       & $5=A_{1},\ 7=B_{1}$ \\
		3  & $(11,13)$     & $11=C_{1},\ 13=E_{1}$ \\
		4  & $(17,19)$     & $17=A_{2},\ 19=B_{2}$ \\
		5  & $(29,31)$     & $29=A_{3},\ 31=B_{3}$ \\
		6  & $(41,43)$     & $41=A_{4},\ 43=B_{4}$ \\
		7  & $(59,61)$     & $59=C_{4},\ 61=E_{3}$ \\
		8  & $(71,73)$     & $71=C_{5},\ 73=E_{4}$ \\
		9  & $(101,103)$   & $101=A_{7},\ 103=B_{7}$ \\
		10 & $(107,109)$   & $107=C_{7},\ 109=E_{6}$ \\
		11 & $(137,139)$   & $137=A_{9},\ 139=B_{9}$ \\
		12 & $(149,151)$   & $149=A_{10},\ 151=B_{10}$ \\
		
	\end{longtable}
	
	\section{Tabelle der ersten Drillinge}
	
	Hier sind die engen Primzahldrillinge der Formen
	\[
	(p,p+2,p+6)\qquad\text{oder}\qquad (p,p+4,p+6),
	\]
	einschließlich des Sonderfalls \((3,5,7)\).
	
	\begin{longtable}{>{\raggedleft\arraybackslash}p{1.2cm} >{\centering\arraybackslash}p{4.7cm} p{7.6cm}}
		\toprule
		Nr. & Drilling & Familienindizes \\
		\midrule
		\endfirsthead
		\toprule
		Nr. & Drilling & Familienindizes \\
		\midrule
		\endhead
		\bottomrule
		\endfoot
		
		1  & $(3,5,7)$       & $3=\text{Sonderfall},\ 5=A_{1},\ 7=B_{1}$ \\
		2  & $(5,7,11)$      & $5=A_{1},\ 7=B_{1},\ 11=C_{1}$ \\
		3  & $(7,11,13)$     & $7=B_{1},\ 11=C_{1},\ 13=E_{1}$ \\
		4  & $(11,13,17)$    & $11=C_{1},\ 13=E_{1},\ 17=A_{2}$ \\
		5  & $(13,17,19)$    & $13=E_{1},\ 17=A_{2},\ 19=B_{2}$ \\
		6  & $(17,19,23)$    & $17=A_{2},\ 19=B_{2},\ 23=C_{2}$ \\
		7  & $(37,41,43)$    & $37=E_{2},\ 41=A_{4},\ 43=B_{4}$ \\
		8  & $(41,43,47)$    & $41=A_{4},\ 43=B_{4},\ 47=C_{3}$ \\
		9  & $(67,71,73)$    & $67=B_{5},\ 71=C_{5},\ 73=E_{4}$ \\
		10 & $(97,101,103)$  & $97=E_{5},\ 101=A_{7},\ 103=B_{7}$ \\
		11 & $(101,103,107)$ & $101=A_{7},\ 103=B_{7},\ 107=C_{7}$ \\
		12 & $(103,107,109)$ & $103=B_{7},\ 107=C_{7},\ 109=E_{6}$ \\
		
	\end{longtable}
	
	\section{Tabelle der ersten Vierlinge}
	
	Wir betrachten Primzahlvierlinge der Form
	\[
	(p,p+2,p+6,p+8).
	\]
	
	\begin{longtable}{>{\raggedleft\arraybackslash}p{1.2cm} >{\centering\arraybackslash}p{5.2cm} p{7.0cm}}
		\toprule
		Nr. & Vierling & Familienindizes \\
		\midrule
		\endfirsthead
		\toprule
		Nr. & Vierling & Familienindizes \\
		\midrule
		\endhead
		\bottomrule
		\endfoot
		
		1  & $(5,7,11,13)$            & $5=A_{1},\ 7=B_{1},\ 11=C_{1},\ 13=E_{1}$ \\
		2  & $(11,13,17,19)$          & $11=C_{1},\ 13=E_{1},\ 17=A_{2},\ 19=B_{2}$ \\
		3  & $(101,103,107,109)$      & $101=A_{7},\ 103=B_{7},\ 107=C_{7},\ 109=E_{6}$ \\
		4  & $(191,193,197,199)$      & $191=C_{11},\ 193=E_{9},\ 197=A_{12},\ 199=B_{12}$ \\
		5  & $(821,823,827,829)$      & $821=A_{38},\ 823=B_{38},\ 827=C_{35},\ 829=E_{32}$ \\
		6  & $(1481,1483,1487,1489)$  & $1481=A_{61},\ 1483=B_{61},\ 1487=C_{59},\ 1489=E_{54}$ \\
		7  & $(1871,1873,1877,1879)$  & $1871=C_{73},\ 1873=E_{68},\ 1877=A_{71},\ 1879=B_{75}$ \\
		8  & $(2081,2083,2087,2089)$  & $2081=A_{79},\ 2083=B_{80},\ 2087=C_{81},\ 2089=E_{74}$ \\
		9  & $(3251,3253,3257,3259)$  & $3251=C_{117},\ 3253=E_{110},\ 3257=A_{118},\ 3259=B_{114}$ \\
		10 & $(3461,3463,3467,3469)$  & $3461=A_{123},\ 3463=B_{121},\ 3467=C_{124},\ 3469=E_{117}$ \\
		
	\end{longtable}
	
	\section{Interpretation}
	
	Diese Tabellen machen zwei Dinge gleichzeitig sichtbar:
	
	\begin{enumerate}
		\item die \emph{Familienform} einer Konstellation, also welche der Klassen \(E,A,B,C\) überhaupt besetzt sind,
		\item den \emph{intrafamiliären Rang} der jeweiligen Primzahl innerhalb ihrer Heimatfamilie.
	\end{enumerate}
	
	Damit erhält jede Primzahl in der Konstellation nicht nur eine Familienzugehörigkeit, sondern auch eine zweite Koordinate:
	\[
	q \longmapsto (F,i),
	\]
	wobei \(F\in\{E,A,B,C\}\) die Heimatfamilie und \(i\) der Familienindex ist.
	
	Beispielsweise wird
	\[
	197 \longmapsto (A,12),\qquad 193 \longmapsto (E,9).
	\]
	
	\section{Strukturelle Beobachtung}
	
	Die Tabellen illustrieren erneut die abgestufte Besetzung des EABC-Raums:
	
	\begin{itemize}
		\item Zwillinge besetzen genau zwei Familien,
		\item Drillinge besetzen genau drei Familien,
		\item Vierlinge besetzen alle vier Familien.
	\end{itemize}
	
	Die Familienindizes zeigen darüber hinaus, dass die Konstellationen zwar die Familienstruktur geordnet durchlaufen, die Zählstände innerhalb der Familien aber nicht notwendig synchron verlaufen.
	
\end{document}